% !TeX root = ../教材题库_全部条目_转写工作稿.tex
% Source: XJ-B1-C05-TH-P168-Y0290
% 第5章 / 5.2 任意角的三角函数 / 教材 p.168 / PDF 第175页
% OCR未校对：公式与文字均可能有误；未完成数学TeX转写；禁止直接发给学生
% 原题及相邻原解截图（核对用）：../题目原图/XJ-B1-C05-TH-P168-Y0290-01.webp
% 本文件只在转写完成后作为题干使用；OCR 条目尚不是可靠题干。
\begin{bankitem}{p.168 思考（含紧邻说明）}{第5章 / 5.2 任意角的三角函数 / 教材 p.168 / PDF 第175页}{XJ-B1-C05-TH-P168-Y0290}{OCR未校对：公式与文字均可能有误；未完成数学TeX转写；禁止直接发给学生}
\begin{ocrdraft}
[教材p.168 / PDF第175页]
\par
？思考前面定义了正弦、余弦、正切这三种三角函数．从定义中可以看出这
\par
些函数是相互关联的，那么，同一个角的三个三角函数值之间有什么关系呢？
\par
如图5．2一7，设仅一Z《rOM是任意角．以点0为圆心作
\par
单位圆与角的终边交于点P，并作角的正弦线DP和余弦
\par
线OD.
\par
在RtAOPD中，由勾股定理得
\par
》DPI2十@02一区丿02
\par
将DP=s1na，OD=cos代人得
\par
图5．2
\par
一7
\par
Sina\ocrglyph{00AF}+\ocrglyph{00AF}COS2；=1．
\par
又由角的终边OP上点P的坐标(cosa，sin及正切函
\par
数的定义得：当co丷乒0时，有
\par
S1n
\par
tan=
\par
COS
\par
这就得到了角的正弦、余弦、正切之间的基本关系式：
\par
Slna\ocrglyph{00AF}f\ocrglyph{00AF}COS：=1，．一
\par
Slna
\par
COS
\end{ocrdraft}
\end{bankitem}
